% ej_coord_R2_04.tex
%
% Copyright (C) 2022--2026 José A. Navarro Ramón <janr.devel@gmail.com>
% Licencia del código GPLv2
% Licencia Creative Commons Recognition Non-Commercial Share-alike.
% (CC-BY-NC-SA)

\section{Ejercicio 4}
Un campo escalar está definido en coordenadas cartesianas ortogonales en $\symbb{R}^2$
por
\[
  \phi(x,y) = xy + 2x
\]
Nos interesa utilizar un sistema de coordenadas cartesianas oblicuas, cuyos ejes forman
\qty{60}{\degree} entre sí, donde el eje de abcisas de ambos sistemas coinciden.\\
Calcula lo que se pide, primero en general para cualquier sistema oblicuo, y después
para nuestro sistema oblicuo concreto. Por último, realiza el cálculo particularizando para el punto $P(x,y) = (3,2)$, en los casos en los que sea pertinente.\\
a) El gradiente en cartesianas ortogonales y su cuadrado.\\
b) La matriz de cambio de base $\mmm{M}$ y su inversa.\\
c) Las componentes contravariantes en función de las ortogonales.\\
d) Las componentes ortogonales en función de las contravariantes.\\
e) La base covariante en función de la base ortogonal y viceversa.\\
f) La métrica en coordenadas oblicuas y su inversa.\\
g) Las componentes covariantes del gradiente.\\
h) Las componentes contravariantes del gradiente.\\
i) El cuadrado del gradiente en coordenadas oblicuas.\\
j) El campo escalar en función de las componentes contravariantes.\\

\noindent
\textbf{Solución}

\noindent
a) Gradiente en cartesianas ortogonales y su cuadrado.

Derivadas parciales
\[
  \dfrac{\partial\phi}{\partial x} = y + 2
  \,;\hspace{1em}
  \dfrac{\partial\phi}{\partial y} = x
\]

Gradiente y su cuadrado
\begin{align*}
  &\vvv{\nabla}\phi
    = \dfrac{\partial\phi}{\partial x}\,\uvec{\i}
    + \dfrac{\partial\phi}{\partial y}\,\uvec{\j}
    = (y+2)\,\uvec{\i} + x\,\uvec{\j}\\
  &|\vvv{\nabla}\phi|^2
    = (y+2)^2 + x^2
\end{align*}
 
Gradiente y su cuadrado en el punto $x=3$ e $y=2$
\begin{align*}
  &\vvv{\nabla}\phi(3,2)
    = (y+2)\,\uvec{\i} + x\,\uvec{\j}
    = (2+2)\,\uvec{\i} + 3\,\uvec{\j}
    = 4\,\uvec{\i} + 3\,\uvec{\j}\\
  & |\vvv{\nabla}\phi|^2
    = 4^2 + 3^2
    = 25
\end{align*}

\noindent
b) Matriz de cambio de base y su inversa.

Mostramos el cambio de base, según la figura \ref{fig:R2-coord_oblicuas_e1e2}, y
su expresión matricial
\[
  \begin{cases}
    \begin{array}{l}
      \xhat{e}_1 = \uvec{\i}\\
      \xhat{e}_2 = \cos\alpha\,\uvec{\i} + \sin\alpha\,\uvec{\j}
    \end{array}
  \end{cases}
  \longrightarrow\hspace{1em}
  \begin{pmatrix}
    \xhat{e}_1 & \xhat{e}_2
  \end{pmatrix}
  = \begin{pmatrix}
    \uvec{\i} & \uvec{\j}
  \end{pmatrix}
  \,\begin{pmatrix}
    1 & \cos\alpha\\
    0 & \sin\alpha
  \end{pmatrix}
\]

De aquí, se deduce que la matriz de cambio de base es
\[
  \mmm{M}
  =\begin{pmatrix}
    1 & \cos\alpha\\
    0 & \sin\alpha
  \end{pmatrix}
\]

La inversa resulta ser
\begin{align*}
  \mmm{M}^{-1}
  &=
    \dfrac{1}{\det(\mmm{M})}
    \,\text{adj}(\mmm{M}^\transp)
    = \dfrac{1}{\sin\alpha}
    \begin{pmatrix}
      \sin\alpha & -\cos\alpha\\
      0 & 1
    \end{pmatrix}
    = \begin{pmatrix}
      1 & -\cot\alpha\\
      0 & \phantom{-}\csc\alpha
    \end{pmatrix}
\end{align*}

En nuestro sistema oblicuo, con $\alpha = \qty{60}{\degree}$
\begin{align*}
  \mmm{M}
  &=
    \begin{pmatrix}
      1 & \cos\pi/3\\
      0 & \sin\pi/3
    \end{pmatrix}
    = \begin{pmatrix}
      1 & 1/2\\
      0 & \sqrt{3}/2
    \end{pmatrix}\\
  \mmm{M}^{-1}
  &=
    \begin{pmatrix}
      1 & -\cot\pi/3\\
      0 & \csc\pi/3
    \end{pmatrix}
    = \begin{pmatrix}
      1 & -\sqrt{3}/3\\
      0 & 2\sqrt{3}/3
    \end{pmatrix}
\end{align*}

Resumen:
\[
  \text{Matriz de cambio de base general e inversa }
  \begin{cases}
    \begin{array}{l}
      \mmm{M}
      = \begin{pmatrix}
        1 & \cos\alpha\\
        0 & \sin\alpha
      \end{pmatrix}\\
      \mmm{M}^{-1}
      = \begin{pmatrix}
        1 & -\cot\alpha\\
        0 & \csc\alpha
      \end{pmatrix}      
    \end{array}
  \end{cases}
\]

\[
  \text{Matriz de cambio de base $\alpha=\qty{60}{\degree}$ e inversa }
  \begin{cases}
    \begin{array}{l}
      \mmm{M}
      = \begin{pmatrix}
        1 & 1/2\\
        0 & \sqrt{3}/2
      \end{pmatrix}\\
      \mmm{M}^{-1}
      = \begin{pmatrix}
        1 & -\sqrt{3}/3\\
        0 & 2\sqrt{3}/3
      \end{pmatrix}      
    \end{array}
  \end{cases}
\]

\noindent
c) Componentes contravariantes en función de las ortogonales.

Valores de $x^1$ y $x^2$ para el punto genérico $P(x,y)$ en coordenadas ortogonales
\[
  \begin{pmatrix}
    x^1\\
    x^2
  \end{pmatrix}
  =
  \begin{pmatrix}
    1 & -\cot\alpha\\
    0 & \csc\alpha
  \end{pmatrix}
  \,\begin{pmatrix}
    x\\
    y
  \end{pmatrix}
  \longrightarrow
  \begin{cases}
    \begin{array}{l}
      x^1 = x-y\cot\alpha\\
      x^2 = y\csc\alpha
    \end{array}
  \end{cases}
\]
\[
  \begin{array}{ll}
  \text{Para nuestro sistema oblicuo y } P(x,y)&
  \begin{cases}
    \begin{array}{l}
      x^1 = x-\dfrac{\sqrt{3}}{3}\,y\\
      x^2 = \dfrac{2\sqrt{3}}{3}\,y
    \end{array}
  \end{cases}\\
  \text{Para nuestro sistema oblicuo y } P(3,2)&
  \begin{cases}
    \begin{array}{l}
      x^1 = 3-\dfrac{2\sqrt{3}}{3}\\
      x^2 = \dfrac{4\sqrt{3}}{3}
    \end{array}
  \end{cases}
  \end{array}
\]

\noindent
d) Las componentes ortogonales en función de las contravariantes.\\
\[
  \begin{pmatrix}
    x\\
    y
  \end{pmatrix}
  =
  \begin{pmatrix}
    1 & \cos\alpha\\
    0 & \sin\alpha
  \end{pmatrix}
  \,\begin{pmatrix}
    x^1\\
    x^2
  \end{pmatrix}
  \longrightarrow
  \begin{cases}
    \begin{array}{l}
      x = x^1 + x^2\cos\alpha\\
      y = x^2\sin\alpha
    \end{array}
  \end{cases}
\]

\[
  \text{En nuestro sistema oblicuo \qty{60}{\degree} }
  \begin{cases}
    \begin{array}{l}
      x = x^1 + \dfrac{1}{2} x^2\\
      y = \dfrac{\sqrt{3}}{2} x^2
    \end{array}
  \end{cases}
\]

\noindent
e) Transformación de la base ortogonal a la covariante del sistema oblicuo y su inversa.

Con la matriz de cambio de base $\mmm{M}$, se obtienen los vectores de la base
covariante
\[
  \begin{pmatrix}
    \xhat{e}_1 & \xhat{e}_2
  \end{pmatrix}
  =\begin{pmatrix}
    \uvec{\i} & \uvec{\j}
  \end{pmatrix}
  \,\begin{pmatrix}
    1 & \cos\alpha\\
    0 & \sin\alpha
  \end{pmatrix}
  \longrightarrow
  \begin{cases}
    \begin{array}{l}
      \xhat{e}_1 = \uvec{\i}\\
      \xhat{e}_2 = \cos\alpha\,\uvec{\i} + \sin\alpha\,\uvec{\j}
    \end{array}
  \end{cases}
\]
Estos vectores miden las componentes contravariantes de un vector en $\symbb{R}^2$ en
coordenadas oblicuas.

La transformación inversa es
\[
  \begin{pmatrix}
    \uvec{\i} & \uvec{\j}                 
  \end{pmatrix}
  =\begin{pmatrix}
    \xhat{e}_1 & \xhat{e}_2
  \end{pmatrix}
  \,\begin{pmatrix}
    1 & -\cot\alpha\\
    0 & \csc\alpha
  \end{pmatrix}
  \longrightarrow
  \begin{cases}
    \begin{array}{l}
      \uvec{\i} = \xhat{e}_1\\
      \uvec{\j} = -\cot\alpha\,\xhat{e}_1 + \csc\alpha\,\xhat{e}_2
    \end{array}
  \end{cases}
\]
Resumen: En nuestro sistema oblicuo $\alpha=\qty{60}{\degree}$, tenemos
\[
  \begin{array}{ll}
  \text{Transformación de base ortogonal a oblicua }&
  \begin{cases}
    \begin{array}{l}
      \xhat{e}_1 = \uvec{\i}\\
      \xhat{e}_2 = \dfrac{1}{2}\,\uvec{\i} + \dfrac{\sqrt{3}}{2}\,\uvec{\j}
    \end{array}
  \end{cases}\\[4ex]
  \text{Transformación de base covariante a ortogonal }&
  \begin{cases}
    \begin{array}{l}
      \uvec{\i} = \xhat{e}_1\\
      \uvec{\j} = -\dfrac{\sqrt{3}}{3}\,\xhat{e}_1 + \dfrac{2\sqrt{3}}{3}\,\xhat{e}_2
    \end{array}
  \end{cases}
  \end{array}
\]

\noindent
f) Métrica y su inversa.

La métrica se calcula mediante el producto escalar de los vectores de la base covariante
\begin{align*}
  \mmm{g}_{ij}
  &=
    \begin{pmatrix}
      g_{11} & g_{12}\\
      g_{21} & g_{22}
    \end{pmatrix}
    = \begin{pmatrix}
      \xhat{e}_1\cdot \xhat{e}_1 & \xhat{e}_1\cdot\xhat{e}_2\\
      \xhat{e}_2\cdot \xhat{e}_1 & \xhat{e}_2\cdot\xhat{e}_2
    \end{pmatrix}
    = \begin{pmatrix}
      1 & \cos\alpha\\
      \cos\alpha & 1
    \end{pmatrix}
\end{align*}

La métrica inversa es
\begin{align*}
  \mmm{g}^{ij}
  &=
    \dfrac{1}{\det(g_{ij})}\,
    \text{adj}\!\left(g_{ij}^\transp\right)
    = \dfrac{1}{\sin^2\alpha}
    \begin{pmatrix}
      1 & -\cos\alpha\\
      -\cos\alpha & 1
    \end{pmatrix}
\end{align*}

\noindent
Resumen:
\[
  \hspace*{-5.5em}
  \text{Métrica general y su inversa }
  \begin{cases}
    \begin{array}{l}
      \mmm{g}_{ij}
      = \begin{pmatrix}
        1 & \cos\alpha\\
        \cos\alpha & 1
      \end{pmatrix}\\
      \mmm{g}^{ij}
      = \dfrac{1}{\sin^2\alpha}
      \begin{pmatrix}
        1 & -\cos\alpha\\
        -\cos\alpha & 1
      \end{pmatrix}
    \end{array}
  \end{cases}
\]
\[
  \text{Métrica y su inversa para \qty{60}{\degree} }
  \begin{cases}
    \begin{array}{l}
      \mmm{g}_{ij}
      = \begin{pmatrix}
        1 & 1/2\\
        1/2 & 1
      \end{pmatrix}\\
      \mmm{g}^{ij}
      =
      \dfrac{1}{3/4}
      \begin{pmatrix}
        1 & -1/2\\
        -1/2 & 1
      \end{pmatrix}
      = \begin{pmatrix}
        4/3 & -2/3\\
        -2/3 & 4/3
      \end{pmatrix}
    \end{array}
  \end{cases}
\]

\noindent
g) Componentes covariantes del gradiente.

Del apartado \emph{d)}, sabemos las coordenadas ortogonales a partir de las
contravariantes. Las sustituimos en las derivadas parciales en cartesianas ortogonales
de \emph{a)}
\[
  \begin{cases}
    \begin{array}{l}
      x = x^1 + x^2\,\cos\alpha\\
      y = x^2\sin\alpha
    \end{array}
  \end{cases}
\]

Expresamos las componentes ortogonales en función de las contravariantes
\begin{align*}
  \dfrac{\partial\phi}{\partial x}
  &=
    y + 2
    = x^2\sin\alpha + 2\\
  \dfrac{\partial\phi}{\partial y}
  &=
    x
    = x^1 + x^2\cos\alpha
\end{align*}

Calculamos las componentes covariantes del gradiente, teniendo en cuenta las derivadas
parciales de $x$ e $y$ respecto de $x^1$ y $x^2$, que dejamos como ejercicio
\begin{align*}
  \hspace{-14em}
  \nabla\phi_1
  &=
    \dfrac{\partial\phi}{\partial x^1}
    = \dfrac{\partial\phi}{\partial x}\dfrac{\partial x}{\partial x^1}
    + \dfrac{\partial\phi}{\partial y}\dfrac{\partial y}{\partial x^1}
    = (x^2\sin\alpha+2)\cdot 1 + (x^1 + x^2\cos\alpha)\cdot 0\\
  &=
    x^2\sin\alpha + 2
\end{align*}
\begin{align*}
  \hspace{1.3em}
  \nabla\phi_2
  &=
    \dfrac{\partial\phi}{\partial x^2}
    = \dfrac{\partial\phi}{\partial x}\dfrac{\partial x}{\partial x^2}
    + \dfrac{\partial\phi}{\partial y}\dfrac{\partial y}{\partial x^2}
    = (x^2\sin\alpha+2)\cdot\cos\alpha + (x^1 + x^2\cos\alpha)\cdot\sin\alpha\\
  &=
    x^1\sin\alpha + 2x^2\sin\alpha\cos\alpha + 2\cos\alpha
\end{align*}

En nuestro sistema oblicuo \qty{60}{\degree}
\begin{align*}
  \nabla\phi_1
  &=
    x^2\sin\pi/3 + 2
    = \dfrac{\sqrt{3}}{2} x^2 + 2\\
  \nabla\phi_2
  &=
    x^1\sin\pi/3 + x^2\,2\sin\pi/3\cos\pi/3 + 2\cos\pi/3
    = \dfrac{\sqrt{3}}{2}x^1 + \dfrac{\sqrt{3}}{2} x^2 + 1
\end{align*}

Para el sistema oblicuo y el punto $P$
\begin{align*}
  \hspace{-7em}
  \nabla\phi_1
  &=
    \dfrac{\sqrt{3}}{2}\cdot \dfrac{4\sqrt{3}}{3} + 2
    = 2 + 2
    = 4
\end{align*}
\begin{align*}
  \nabla\phi_2
  &=
    \dfrac{\sqrt{3}}{2}\cdot \left(3-2\dfrac{\sqrt{3}}{3}\right)
    + \dfrac{\sqrt{3}}{2}\cdot\dfrac{4\sqrt{3}}{3} + 1
    = \dfrac{3\sqrt{3}}{2} - 1 + 2 + 1\\
  &=
    2+\dfrac{3\sqrt{3}}{2}
\end{align*}

\noindent
Resumen:
\[
  \hspace{1em}
  \text{Sistema oblicuo general }
  \begin{cases}
    \begin{array}{l}
      \nabla\phi_1 = x^2\sin\alpha + 2\\
      \nabla\phi_2 = x^1\sin\alpha + 2x^2\sin\alpha\cos\alpha + 2 \cos\alpha
    \end{array}
  \end{cases}
\]
\[
  \begin{array}{ll}
    \text{Sistema oblicuo \qty{60}{\degree} } &
    \begin{cases}
      \begin{array}{l}
        \nabla\phi_1 = \dfrac{\sqrt{3}}{2} x^2 + 2\\[2ex]
        \nabla\phi_2 =\dfrac{\sqrt{3}}{2} x^1 + \dfrac{\sqrt{3}}{2} x^2 + 1
      \end{array}
    \end{cases}\\[2ex]
    \text{Sistema oblicuo y punto } P(3,2) &
    \begin{cases}
      \begin{array}{l}
        \nabla\phi_1 = 4\\
        \nabla\phi_2 = 2 + \dfrac{3\sqrt{3}}{2}
      \end{array}
    \end{cases}
  \end{array}
\]

\noindent
h) Componentes contravariantes del gradiente.

\vspace*{1ex}
Se calculan mediante $\nabla\phi^i = g^{ij} \nabla\phi_j$
\begin{align*}
  \nabla\phi^1
  &=
    g^{11} \nabla\phi_1 + g^{12}\nabla\phi_2\\
  &=
    \dfrac{1}{\sin^2\alpha} \left(x^2\,\sin\alpha + 2\right)
    - \dfrac{\cos\alpha}{\sin^2\alpha}
    \,\left(x^1\sin\alpha
    + 2x^2\sin\alpha\cos\alpha
    + 2\cos\alpha\right)\\
  &=
    - x^1\,\cot\alpha
    + x^2\,\dfrac{1-2\cos^2\alpha}{\sin\alpha}
    + 2
\end{align*}
\begin{align*}
  \nabla\phi^2
  &=
    g^{21} \nabla\phi_1 + g^{22}\nabla\phi_2\\
  &=
    -\dfrac{\cos\alpha}{\sin^2\alpha} \left(x^2\,\sin\alpha + 2\right)
    + \dfrac{1}{\sin^2\alpha}
    \,\left(x^1\sin\alpha
    + x^2\,2\sin\alpha\cos\alpha
    + 2\cos\alpha\right)\\
  &=
    x^1\csc\alpha
    + x^2\left(%
    2\,\cot\alpha
    -\cot\alpha
    \right)
    - \cancelout{\dfrac{2\cos\alpha}{\sin^2\alpha}}
    + \cancelout{\dfrac{2\cos\alpha}{\sin^2\alpha}}
    = x^1\csc\alpha
    + x^2\cot\alpha
\end{align*}

Para el sistema oblicuo con $\alpha= \qty{60}{\degree}$
\begin{align*}
  \nabla\phi^1
  &=
    -x^1 \cot{\pi/3}
    + x^2\,\dfrac{1-2\cos^2\pi/3}{\sin\pi/3}
    + 2\\
  &=
    -\dfrac{\sqrt{3}}{3}x^1
    +  \dfrac{1-2\,(1/4)}{\sqrt{3}/2} x^2
    + \dfrac{2\cos\pi/3}{\sin^2\pi/3}
    = -\dfrac{\sqrt{3}}{3}x^1
    + \dfrac{\sqrt{3}}{3}x^2
    + 2\\
  \nabla\phi^2
  &=
    x^1 \csc\pi/3
    + x^2 \cot\pi/3
    = \dfrac{2\sqrt{3}}{3} x^1 + \dfrac{\sqrt{3}}{3} x^2
\end{align*}

Aplicado al punto $P(3,2)$
\begin{align*}
  \nabla\phi^1
  &=
    -\dfrac{\sqrt{3}}{3}\left(%
    3-\dfrac{2\sqrt{3}}{3}
    \right)
    + \dfrac{\sqrt{3}}{3}\, \dfrac{4\sqrt{3}}{3}
    + 2
    = -\sqrt{3} + \dfrac{2}{3} + \dfrac{4}{3} + 2
    = 4 - \sqrt{3}\\
  \nabla\phi^2
  &=
    \dfrac{2\sqrt{3}}{3} x^1 + \dfrac{\sqrt{3}}{3} x^2
    = \dfrac{2\sqrt{3}}{3} \left(3-\dfrac{2\sqrt{3}}{3}\right)
    + \dfrac{\sqrt{3}}{3}\,\dfrac{4\sqrt{3}}{3}\\
  &=
    2\sqrt{3}
    - \dfrac{4}{3}
    + \dfrac{4}{3}
    = 2\sqrt{3}
\end{align*}

\noindent
Resumen:
\[
  \text{Sistema oblicuo general }
  \hspace{1.2em}
  \begin{cases}
    \begin{array}{l}
      \nabla\phi^1 = -x^1\cot\alpha
      + x^2\,\dfrac{1-2\cos^2\alpha}{\sin\alpha}
      + 2\\
      \nabla\phi^2 = x^1\csc\alpha + x^2 \cot\alpha
    \end{array}
  \end{cases}
\]
\[
  \begin{array}{ll}
    \text{Sistema oblicuo \qty{60}{\degree} } &
    \begin{cases}
      \begin{array}{l}
        \nabla\phi^1 = -\dfrac{\sqrt{3}}{3} x^1 + \dfrac{\sqrt{3}}{3} x^2 + 2\\[2ex]
        \nabla\phi^2 =\dfrac{2\sqrt{3}}{3} x^1 + \dfrac{\sqrt{3}}{3} x^2
      \end{array}
    \end{cases}\\[2ex]
    \text{Sistema oblicuo y punto } P &
    \begin{cases}
      \begin{array}{l}
        \nabla\phi^1 = 4 - \sqrt{3}\\
        \nabla\phi^2 = 2\sqrt{3}
      \end{array}
    \end{cases}
  \end{array}
\]

\noindent
i) Cuadrado del gradiente en coordenadas oblicuas.\\

Realizamos el cálculo por partes
\begin{align*}
    \nabla\phi^1\nabla\phi_1
  &=
    \left[%
    -x^1\cot\alpha
    + x^2\,\dfrac{1-2\cos^2\alpha}{\sin\alpha}
    + 2
    \right]
    \left(%
    x^2\sin\alpha + 2
    \right)\\
  &\!=
    \!- 2x^1\!\cot\alpha
    +\!\left(x^2\right)^2\!\left(1-2\cos^2\alpha\right)
    \!+ 2x^2\!\left(\!2\sin\alpha - \dfrac{\cos^2\alpha}{\sin\alpha}\right)\!
    -x^1x^2\cos\alpha
    + 4
\end{align*}
\begin{align*}
  \hspace*{-5em}
  \nabla\phi^2\nabla\phi_2
  &=
    \left(x^1\csc\alpha + x^2 \cot\alpha\right)
    \left(x^1\sin\alpha + x^2 2\sin\alpha\cos\alpha + 2 \cos\alpha\right)\\
  &=
    \left(x^1\right)^2
    + 2x^1\cot\alpha
    + 2\left(x^2\right)^2\cos^2\alpha
    + 2x^2\dfrac{\cos^2\alpha}{\sin\alpha}
    + 3x^1x^2\cos\alpha
\end{align*}
\begin{align*}
  |\vvv{\nabla}\phi|^2
  &=
    \nabla\phi^1\nabla\phi_1 + \nabla\phi^2\nabla\phi_2\\
  &=
    \left(x^1\right)^2
    + \left(x^2\right)^2
    + 4x^2\sin\alpha
    + 2x^1x^2\cos\alpha
    + 4
\end{align*}

Para nuestro sistema oblicuo
\begin{align*}
  |\vvv{\nabla}\phi|^2
  &=
    \left(x^1\right)^2
    + \left(x^2\right)^2
    + 4x^2\sin\alpha
    + 2x^1x^2\cos\alpha
    + 4\\
  &=
    \left(x^1\right)^2
    + \left(x^2\right)^2
    + 4x^2\,\dfrac{\sqrt{3}}{2}
    + 2x^1x^2\,\dfrac{1}{2}
    + 4\\
  &=
    \left(x^1\right)^2
    + \left(x^2\right)^2
    + 2\sqrt{3} x^2
    + x^1 x^2
    + 4
\end{align*}

Cuando se aplica al punto $P$
\begin{align*}
  |\vvv{\nabla}\phi|^2
  &=
    \left(x^1\right)^2
    + \left(x^2\right)^2
    + 2\sqrt{3}\, x^2
    + x^1 x^2
    + 4\\
  &=
    \left(3-\dfrac{2\sqrt{3}}{3}\right)^2
    + \left(\dfrac{4\sqrt{3}}{3}\right)^2
    + 2\sqrt{3}\,\dfrac{4\sqrt{3}}{3}
    + \left(3-\dfrac{2\sqrt{3}}{3}\right)\,\dfrac{4\sqrt{3}}{3}
    + 4\\
  &=
    9
    - \cancelout{4\sqrt{3}}
    + \dfrac{4}{3}
    + \dfrac{16}{3}
    + 8
    + \cancelout{4\sqrt{3}}
    - \dfrac{8}{3}
    + 4
    = 21 + 4
    = 25
\end{align*}

\noindent
Resumen:
\[
  |\vvv{\nabla}\phi|^2 =
  \begin{cases}
    \begin{array}{l}
      \text{Sistema oblicuo general: }
      \left(x^1\right)^2 + \left(x^2\right)^2 + 4x^2\sin\alpha + 2x^1x^2\cos\alpha + 4\\
      \text{Sistema oblicuo \qty{60}{\degree}}:
      \left(x^1\right)^2 + \left(x^2\right)^2 + 2\sqrt{3} x^2 + x^1 x^2 + 4\\
      \text{Sistema oblicuo \qty{60}{\degree} y punto } $P$ \text{ : }
      25
    \end{array}
  \end{cases}
\]

\noindent
j) Campo escalar en función de las componentes contravariantes.

\begin{align*}
  \phi
  &=
    xy+2x
    = \left(x^1+x^2\cos\alpha\right) x^2\sin\alpha + 2 \left(x^1+x^2\cos\alpha\right)\\
  &=
    x^1x^2\sin\alpha
    + \left(x^2\right)^2\sin\alpha\cos\alpha
    + 2 x^1
    + 2x^2\cos\alpha\\
  &=
    2 x^1
    + \left(x^2\right)^2\sin\alpha\cos\alpha
    + 2x^2\cos\alpha
    + x^1x^2\sin\alpha
\end{align*}

En nuestro sistema oblicuo concreto
\begin{align*}
  \phi
  &=
    2 x^1
    + \left(x^2\right)^2\sin\alpha\cos\alpha
    + 2x^2\cos\alpha
    + x^1x^2\sin\alpha\\
  &=
    2x^1
    + \left(x^2\right)^2\dfrac{\sqrt{3}}{2}\,\dfrac{1}{2}
    + 2x^2\dfrac{1}{2}
    + x^1x^2\dfrac{\sqrt{3}}{2}\\
  &=
    2x^1
    + \dfrac{\sqrt{3}}{4}\left(x^2\right)^2
    + x^2
    + \dfrac{\sqrt{3}}{2}x^1x^2
\end{align*}

Su valor para el campo escalar para el punto $(x,y) = (3, 2)$, en coordenadas
ortogonales
\begin{align*}
  \phi
  &=
    2x^1
    + \dfrac{\sqrt{3}}{4}\left(x^2\right)^2
    + x^2
    + \dfrac{\sqrt{3}}{2}x^1x^2\\
  &=
    2\left(3-\dfrac{2\sqrt{3}}{3}\right)
    + \dfrac{\sqrt{3}}{4}\left(\dfrac{4\sqrt{3}}{3}\right)^2
    + \dfrac{4\sqrt{3}}{3}
    + \dfrac{\sqrt{3}}{2} \left(3-\dfrac{2\sqrt{3}}{3}\right)
    \left(\dfrac{4\sqrt{3}}{3}\right)\\
  &=
    6
    - \cancelout{\dfrac{4\sqrt{3}}{3}}
    + \cancelout{\dfrac{4\sqrt{3}}{3}}
    + \cancelout{\dfrac{4\sqrt{3}}{3}}
    + 6
    - \cancelout{\dfrac{4\sqrt{3}}{3}}
    = 12
\end{align*}

En coordenadas ortogonales da
\[
  \phi
  = xy + 2x
  = 3\cdot 2 + 2 \cdot 3
  = 12
\]

\noindent
Resumen:
\[
  \phi =
  \begin{cases}
    \begin{array}{l}
      \text{Sistema oblicuo general: }
      2x^1 + \left(x^2\right)^2\sin\alpha\cos\alpha + 2x^2\cos\alpha + x^1x^2\sin\alpha\\
      \text{Sistema oblicuo \qty{60}{\degree}}:
      2x^1 + \dfrac{\sqrt{3}}{4}\left(x^2\right)^2 + x^2 + \dfrac{\sqrt{3}}{2}x^1x^2\\
      \text{Sistema oblicuo \qty{60}{\degree} en } $P$ \text{ : }
      12
    \end{array}
  \end{cases}
\]









%%% Local Variables:
%%% mode: latex
%%% TeX-engine: luatex
%%% TeX-master: "../retazosmatematicas.tex"
%%% End:
